\chapter*{\rm\bfseries Abstract}
\addcontentsline{toc}{chapter}{Abstract}

Robotic grasping of thin, deformable fabric remains an unsolved problem
in industrial automation. Two competing gripper geometries are routinely
proposed. Rigid parallel-jaw fingers offer precise force control but
narrow operational tolerance, while compliant finray-effect fingers
trade force precision for a wider operational envelope. In this work,
an experimental evaluation of that trade is presented on the task of
picking a single flat cloth off a surface, comparing two rigid
parallel-jaw configurations and four 3D-printed finray configurations,
six gripper variants in all, on a Franka Research 3 arm. Press depth and tilt
angle were chosen as the evaluation axes because they are the placement
errors that a deployed system cannot fully control. A low-accuracy or
coarsely calibrated arm carries millimetre-scale surface-height error
and degree-scale orientation error, and both of them must be absorbed
by the gripper. Stationary contact force was the primary outcome and
grasp success the secondary outcome.

Across 840 analysed trials, the rigid jaws held cloth only within a
$1.5$\,mm depth window, or $2.5$\,mm with a compliant pad, bounded by a
force safety limit. However, every finray variant grasped across the
full $10$\,mm calibrated analysis window at every orientation tested.
In addition, all four finray variants were found to share a single
force-depth curve once each was normalised by its own peak force, with
a mean pairwise error below $1\%$. The finger geometry therefore sets
the scale of the contact force, although it does not change the shape
of the response. The tilt response itself proved to be depth-gated.
Finray contact force is strongly tilt-sensitive at the contact
boundary, however it becomes robust to tilt once the finger is pressed
past the buckling knee, where the force spread across all tested
orientations falls to between $1.3$ and $2.4$\,N against knee peaks of
$31$ to $53$\,N. Therefore, compliant grippers only obtain this
robustness when they are pressed past a threshold depth, and a control
policy that works near the contact boundary for delicate fabric
handling should account for the sensitivity that remains there.
